%------------------------------------------------------------------------------%
\section{Inequação do Segundo Grau}

Para resolver uma inequação do segundo grau, você deve resolver a equação do
segundo grau correspondente e fazer o estudo do sinal.

%------------------------------------------------------------------------------%
\begin{example}
  Resolva a inequação $x^2-5x+6>0$

  \solution
  Inicialmente devemos resolver a equação do segundo grau
  \[
    x^2 - 5x + 6 = 0,
  \]
  cujas raízes, obtidas pela fórmula de Bhaskara são $x=2$ e $x=3$.
  Portanto, temos uma parábola com concavidade para cima ($a>0$)
  e tocando o eixo $x$ em dois pontos.
  Então ``fora da parábola'' colocamos o sinal de $a$ e dentro sinal
  contrário a $a$ como mostra a figura.

  \IncludeGraphicCentredEx{ineq1}[0.5]

  Logo, a solução da inequação é
  \[
    \{ x \in \R \st x < 2 \text{ ou } x > 3 \}.
  \]
\end{example}
%------------------------------------------------------------------------------%

%------------------------------------------------------------------------------%
\begin{example}
  Encontre os valores de $x$ que satisfazem a desigualdade $9 \leq x^2$.

  \solution
  Note que $9 \leq x^2 $ é equivalente a $-x^2+9 \leq 0$.
  Portanto, precisamos resolver a inequação do segundo grau
  \[
    -x^2 + 9 \leq 0.
  \]
  A equação do segundo grau correspondente é
  \[
    -x^2 + 9 = 0,
  \]
  onde $a=-1$, $b=0$ e $c=9$. As raízes dessa equação são $x=-3$ e $x=3$.
  O estudo do final é mostrado na figura.

  \IncludeGraphicCentredEx{ineq2}[0.5]

  Logo, a solução pedida é
  \[
    \{ x \in \R \st x \leq -3 \text{ ou } x \geq 3 \}.
  \]
\end{example}
%------------------------------------------------------------------------------%
